% Speaker companion for milestone3 --- SUTVA statistical-battery slides.
% Per-slide talking points, background context, and Q&A backup material.
% Compile: pdflatex milestone3_speaker_companion.tex

\documentclass[11pt,a4paper]{article}

\usepackage[margin=2cm,top=1.8cm,bottom=2cm]{geometry}
\usepackage[hidelinks]{hyperref}
\usepackage{xcolor}
\usepackage{enumitem}
\usepackage{titlesec}
\usepackage{parskip}
\usepackage{tcolorbox}
\usepackage{amsmath,amssymb}
\usepackage{booktabs}

\definecolor{accent}{HTML}{8B0000}
\definecolor{muted}{HTML}{555555}
\definecolor{lightbg}{HTML}{F7F4F0}

\titleformat{\section}{\large\bfseries\color{accent}}{\thesection.}{0.6em}{}
\titlespacing*{\section}{0pt}{1.2em}{0.2em}
\titleformat{\subsection}{\normalsize\bfseries}{}{0em}{}
\titlespacing*{\subsection}{0pt}{0.6em}{0.1em}

\setlist[itemize]{leftmargin=1.2em,itemsep=0.15em,topsep=0.2em}

\newcommand{\timetag}[1]{\hfill\textcolor{muted}{\small[#1]}}
\newcommand{\highlight}[1]{\textcolor{accent}{\textbf{#1}}}

\title{\bfseries Milestone 3 --- Speaker Companion: SUTVA Statistical Battery}
\author{Per-slide guidance, background, and Q\&A backup}
\date{}

\begin{document}

\maketitle
\vspace{-2em}

\begin{tcolorbox}[colback=lightbg,colframe=accent!40,arc=2pt]
\textbf{One-line spine for these two slides.} \emph{The economic audit is the SUTVA defense; the
permutation test and HLT \emph{confirm} it; the break tests are informative loading-stability
diagnostics, not exclusion rules.} If you remember nothing else, say that sentence --- it is also
exactly what the code now encodes (\texttt{config.BREAKPOINT\_EXCLUDE} is empty; pool membership
follows the audit).
\end{tcolorbox}

% =====================================================================
\section{How we defend donor selection (SUTVA) \timetag{60s}}

\subsection{Talking points}
\begin{itemize}
  \item \highlight{Primary defense is the economic audit} --- the SCM-canonical approach (Abadie 2021
    \S3.1: donor exclusion rests on \emph{substantive} grounds, not statistical pretests). Screen every
    candidate against two contamination channels: the oil-supply path (energy complex, petrocurrencies
    --- categorically out) and the disruption itself (Hormuz routing, Russia/Ukraine trade --- donor by donor).
  \item \highlight{Statistical cross-check --- the two tests on the slides:}
    \begin{itemize}
      \item \highlight{Permutation mean-shift} (Chow 1960 framing; Lehmann \& Romano permutation
        reference) --- the literal SUTVA question (did the donor's own series shift at $T_0$?), Brent-free,
        nonparametric. The validation doc designates this as the minimum load-bearing statistical
        confirmation, and it is the test that produces the lone flag (CNY).
      \item \highlight{HLT trend break} --- Brent-free and I(0)/I(1)-robust, so a break at $T_0$
        \emph{unambiguously} implicates the donor itself. The cleanest of the battery; flags nothing.
    \end{itemize}
  \item \highlight{Bai--Perron and KS are diagnostics, not exclusion rules} --- say this explicitly.
    Bai--Perron regresses on \emph{treated} Brent, so a $T_0$ break is circular (could just be Brent
    moving); it measures \emph{loading stability}, not SUTVA. KS over-rejects under macro-regime change.
    Both are reported but neither auto-excludes a donor.
\end{itemize}

\subsection{The other at-$T_0$ test (run, not on the slides): Wilcoxon event-window}
A third at-$T_0$ test sits in the battery but is kept off the slides:
\begin{itemize}
  \item \highlight{Wilcoxon event-window} (Brown \& Warner 1985, event-study tradition): compares the
    donor's returns in a \emph{short window around} $T_0$ against a control window, rank-based. Where the
    permutation test looks for a \emph{persistent mean shift} across the whole pre/post split, the
    event-window looks for a \emph{localized return reaction} right at the event --- they are complementary.
  \item \highlight{Why it is not shown:} on the retained pool it flags \emph{nothing} --- it does not even
    pick up CNY ($p \approx 0.44$ vs the permutation test's $p \approx 0.002$) --- and it is lower-powered
    on the handful of event-window observations. So it adds no information to the result; the permutation
    test carries the lone flag on its own. Mention it only if asked ``did you run more than one at-$T_0$ test?''
  \item In the combined BH-FDR flag the pipeline computes, the flag is ``permutation \emph{or} event-window'';
    since event-window fires on nothing here, the combined flag equals the permutation test alone.
\end{itemize}

\subsection{Behind the slide}
The hierarchy matters: audit is what a reviewer expects and what carries the argument; everything
statistical is supplementary. We deliberately do \emph{not} let a test exclude a donor automatically ---
the test flags a candidate, the researcher decides. This is the ``report and interpret'' convention of
the SCM literature.

\subsection{If asked}
\begin{itemize}
  \item \textit{Why isn't Bai--Perron a SUTVA test?} It regresses each donor on Brent's returns. Brent
    is the treated unit, so at $T_0$ Brent itself jumps; a break in the donor--Brent relationship there
    cannot be separated from Brent's own move. It tells you whether the donor's \emph{loading} is stable
    across the fitting window (a parallel-fit concern), not whether the donor was treated.
  \item \textit{Why drop KS?} It rejects $\sim$21 of 32 donors on Russia --- but that is the macro-regime
    change (COVID recovery $\to$ 2022 tightening), not per-donor treatment. Our own flag rule excludes it
    for that reason; it is a regime diagnostic, not a SUTVA test.
  \item \textit{Is HLT definitive?} No --- a non-rejection is failure-to-reject, not proof of no
    interference, and HLT only catches \emph{trend} breaks. It is the cleanest \emph{statistical}
    triangulation, but the audit remains primary.
\end{itemize}

% =====================================================================
\section{SUTVA test: no retained donor is treated at $T_0$ \timetag{75s}}

\subsection{The table (verified, 19-donor retained pool)}
\begin{center}
\begin{tabular}{@{}l c c@{}}
  \toprule
  \textbf{Flagged at $T_0$ (of 19 retained)} & \textbf{Hormuz 2026} & \textbf{Russia 2022} \\
  \midrule
  HLT trend-break (Brent-free $\Rightarrow$ SUTVA)   & 0 / 19 & 0 / 19 \\
  Permutation mean-shift \emph{(on slides)}          & 0 / 19 & 1 / 19 (CNY) \\
  Wilcoxon event-window \emph{(run, not shown)}      & 0 / 19 & 0 / 19 \\
  Bai--Perron (loading stability; appendix)          & 0 / 19 & 1 / 19 (VIX) \\
  \bottomrule
\end{tabular}
\end{center}

\subsection{Talking points}
\begin{itemize}
  \item \highlight{HLT flags nothing at either event} --- the headline. Since it is Brent-free, this is
    the strongest clean read: no retained donor's own trend breaks at the event.
  \item \highlight{Two Russia flags, both retained on substantive grounds:}
    \begin{itemize}
      \item \highlight{CNY} (permutation test) --- the shift is the concurrent China zero-COVID
        reopening, not Russia-treatment. The test cannot separate the two; the audit judges it clean.
      \item \highlight{VIX} (Bai--Perron) --- a generic risk-off spike around the invasion, i.e.\
        \emph{intended} common-factor co-movement, which is what SCM exploits, not contamination.
    \end{itemize}
  \item \highlight{Tests inform; the audit decides.} Neither flag triggers an exclusion --- that is the
    policy, and it is why both donors stay in the 19-donor pool.
\end{itemize}

\subsection{Behind the slide}
VIX was \emph{previously} excluded by an automatic ``breakpoint rule.'' We removed that rule after
two checks: (i) the Bai--Perron break does not reproduce when the test is re-run on the analysis window
alone (it sits in the trimming dead zone; the contemporaneous early-2022 break is better attributed to
HYG), and (ii) a VIX-in/out sensitivity moves the ensemble premium by only $-0.2$ pp (Russia) and
$+0.6$ pp (Hormuz). Audit-vs-test agreement overall: Russia 18 / 32, Hormuz 31 / 32.

\subsection{If asked}
\begin{itemize}
  \item \textit{If a test flags VIX, why keep it?} The flag is a loading-stability signal, not SUTVA;
    it is sample-dependent (does not survive an analysis-window re-run); and it is immaterial to the
    estimate ($\leq 0.6$ pp; see \texttt{scripts/vix\_sensitivity.py}). Excluding it would also force
    dropping a clean, high-information donor from Hormuz purely to keep a shared pool.
  \item \textit{Then why show Bai--Perron at all?} Honesty and completeness --- it is a legitimate
    fit-stability diagnostic and it independently re-flags the audit's known cases (e.g.\ Wheat). We
    present it as what it is, not as a SUTVA test.
  \item \textit{Why does the permutation test flag CNY but not the audit?} The test sees ``CNY moved at
    $T_0$''; it cannot attribute that to Russia versus the concurrent China reopening. The audit, which
    reasons about causal channels, places CNY's move on the China-policy channel --- off the oil-supply /
    disruption path --- so it stays.
  \item \textit{Couldn't HLT just be under-powered on a short window?} On the full-history spec it is
    well-powered and clean. A window-restricted HLT does throw a few rejections, but at off-$T_0$ dates
    (COVID, 2025 metals moves) --- volatility episodes, not focal-event treatment --- so the full-history
    result is the right one to read.
\end{itemize}

\end{document}
